\begin{trapbox}
\textbf{误区一：距离比为 $1$ 就不必检查 $F\notin\ell$。}
\begin{center}\begin{tikzpicture}[font=\small,>=stealth]
\draw[black!30,rounded corners=2pt] (0,0) rectangle (14.4,8.2);
\node[anchor=west,font=\small\bfseries] at (.5,7.72) {距离相等还须有前提：$F$ 不在准线上};
\begin{scope}[shift={(3.6,4)},scale=.4]
\draw[violet,thick] (-3,-4)--(-3,4);\draw[blue,very thick,domain=-4:4,samples=121,variable=\t] plot ({\t*\t/12},\t);\node[fill=black,circle,inner sep=1.8pt] at (3,0) {};\node[below right=7pt] at (3,0) {$F$};
\end{scope}
\node at (3.6,6.6) {$F=(3,0),\ \ell:x=-3$};
\node at (3.6,1.5) {等距轨迹为抛物线};
\begin{scope}[shift={(10.8,4)},scale=.4]
\draw[violet,thick] (0,-4)--(0,4);\draw[blue,very thick] (-4,0)--(5,0);\node[fill=black,circle,inner sep=1.8pt] at (0,0) {};\node[below right=7pt] at (0,0) {$F$};
\end{scope}
\node at (10.8,6.6) {$F=(0,0),\ \ell:x=0$};
\node at (10.8,1.5) {等距轨迹为直线};

\node[anchor=west] at (.7,.55) {同尺度；右图 $\sqrt{x^2+y^2}=|x|$ 等价于 $y=0$，在 $P=F$ 处距离比无定义};
\end{tikzpicture}\end{center}

右图等距条件给出整条直线 $y=0$。除去 $P=F$ 后虽然距离比仍为 $1$，轨迹也不是非退化抛物线。

\textbf{误区二：把 $p$ 当作顶点到焦点的距离。}
在 $y^2=2px$ 中，$p$ 是焦点到准线的距离，$VF=p/2$，一次项系数为 $2p$。
标准位置和平移位置必须先分清，再判断焦点、准线。

\textbf{误区三：所有圆锥曲线的焦半径都等于准距。}
等距是统一定义 $e=1$ 的情形；椭圆为 $0<e<1$，双曲线为 $e>1$。
不要把“统一定义的特例”写成“与统一定义不同”。
\end{trapbox}
\clearpage

% ===== 以下为 C041 防错与条件补充 =====
\begin{trapbox}
\textbf{误区四：把 $P$ 连向准线上任何一点的线段都当作准距。}
只允许取垂足 $H$：
\[d(P,\ell)=PH=\min_{A\in\ell}PA.\]
一般准线上点 $A$ 满足 $PA\ge PH$；若随意将 $PF=PA$ 当定义，轨迹会改变。

\medskip
\textbf{误区五：平方之后没有说明为什么是充要条件。}
这里 $PF\ge0$、$d(P,\ell)\ge0$，
所以 $PF=d(P,\ell)\iff PF^2=d(P,\ell)^2$。
但若把\textbf{无符号距离}误写成有符号的 $x+p/2$，就可能出错。
正确顺序：$PH=|x+p/2|$；先由方程推出 $x\ge0$，再去掉绝对值。

\medskip
\textbf{误区六：两种参数习惯混用。}
若写 $y^2=2px$，则 $VF=p/2$、$FD=p$；
若写 $y^2=4ax$，则 $VF=a$、$FD=2a$。两式等价，$p=2a$。
例如 $y^2=12x$ 的焦点是 $(3,0)$，\textbf{不是} $(6,0)$。

\medskip
\textbf{误区七：不看焦点与准线的相对位置，就断言开口。}
$F=(3,0)$、$\ell:x=-3$ 的顶点是 $(0,0)$，向右开口；
同一焦点配 $\ell:x=5$，顶点 $(4,0)$，\textbf{反而向左开口}。
只报焦点而不报准线，既不能确定顶点，也不能确定轨迹。
\end{trapbox}
\clearpage
\begin{trapbox}
\textbf{误区八：已知焦半径 $L$ 就默认有两个对应点。}
以 $y^2=2px$ 为例，$PF=L$ 等价于 $x=L-p/2$，代回可知：
\[
\begin{cases}
L<p/2 &: \text{无点};\\
L=p/2 &: \text{仅顶点一处};\\
L>p/2 &: \text{上下对称的两点}.
\end{cases}
\]
在一般朝向中“两个点”表现为关于\textbf{对称轴}对称，不一定是上下两个。

\medskip
\textbf{误区九：斜准线只靠平移就能变成 $y^2=2px$。}
坐标平移不会改变一条直线的斜率。
对于 $\ell:x+y=0$，须用点线距离公式直接列式，
或转动坐标轴使一条新轴与 $FD$ 平行；不能只改成 $(y-k)^2=2p(x-h)$。

\medskip
\textbf{误区十：把“圆过焦点且切准线”误读成“圆与抛物线相切”。}
本定义里的圆\textbf{经过固定焦点}，\textbf{与固定准线相切}，
它的\textbf{圆心}组成抛物线。没有声称这些圆都与抛物线相切。

\medskip
\textbf{三项快速自检：}
(1) 焦点是否确实不在准线上；
(2) 焦点、顶点与准线垂足是否共线，且其所在直线与准线垂直；
(3) 任选一个简单的候选点，代入“到点距离=到线垂距”是否成立。
\end{trapbox}
